% C.40: root and DE independent full proof/code audits passed.
\subsection{Independent local coordinates before and after descent}

\begin{statement}{C.40}{Proved}{The simultaneous gentle image at every prime at least five}
For a prime $p\geq5$, let $G_{0,p}$ denote the group of actual
normalized pro-$p$ gentle pairs, with conjugator in the pro-$p$
commutator subgroup, equivalently their faithful pure outer classes:
the two cycle identities hold and the induced
map of the free pro-$p$ group is an automorphism. Let
$\lambda_p:G_{0,p}\to\mathbb Z_p^\times$ be its multiplier.
Then the image of the actual full group under
\[
 \widehat{\mathrm{GT}}_0\longrightarrow
 \prod_{p\geq5}(G_{0,p}\times G_{0,p}),\qquad
 g\longmapsto (g_p,D(g)_p)_{p\geq5}
                                                        \tag{C.40.1}
\]
is exactly
\[
 \prod_{p\geq5}
       (G_{0,p}\mathbin\times_{\mathbb Z_p^\times}G_{0,p}).
                                                        \tag{C.40.2}
\]
In particular, the multiplier-one image is the whole product
of the two multiplier-one local gentle groups, simultaneously
at all these primes. This determines the gentle-group image
asked for in Task C.5. It does not determine the simultaneous
image of the full pentagon group, or assert any inclusion of
that group in $C^\varphi$.
\end{statement}

\begin{proof}
We first prescribe multiplier-one elements at a finite set of
primes $p\geq5$. The construction retains both local actions
in one actual quotient of the modular source, then checks its
entire source-centralization. No assertion that a relative
quotient is an absolute pronilpotent quotient is used.

\paragraph{The six actual fillings and their finite base.}
Put
\[
 \Gamma=\widehat{C_2*C_3}=\langle s,r\mid s^2=r^3=1\rangle,
 \quad F=\ker(\Gamma\to S_3),\quad
 x=(rs)^2,\quad y=(sr)^2,\quad z=(xy)^{-1}.
                                                        \tag{C.40.3}
\]
Write $\sigma=\operatorname{Int}(s)$ and
$\rho=\operatorname{Int}(r)$. Direct reduction in the free
product gives the exact formulas
\[
 \sigma(x)=y,\quad\sigma(y)=x,\qquad
 \rho(x)=z,\quad\rho(y)=x.                              \tag{C.40.4}
\]
Thus $\sigma^2=\rho^3=1$ as actual automorphisms of $F$.
Let $\Gamma'$ be a separately marked copy of $\Gamma$, and
let the filling $\kappa:F\twoheadrightarrow\Gamma'$ be
\[
       \kappa(x)=s',\qquad \kappa(y)=s'r',\qquad
                         \kappa(z)=r'^{-1}.
                                                        \tag{C.40.5}
\]
For $k=0,1,2$ use the six maps
$\kappa_k=\kappa\rho^k$ and
$\kappa'_k=\kappa\sigma\rho^k$. These are actual maps,
not a list of unrelated outer isomorphism types. The order-two
automorphism
\[
                     t_k=\rho^{-k}\sigma\rho^k
                                                        \tag{C.40.6}
\]
interchanges $\kappa_k$ and $\kappa'_k$ literally.

In the standard two-dimensional representation of the target
$S_3$, choose
\[
 S=\begin{pmatrix}0&1\\1&0\end{pmatrix},\quad
 R=\begin{pmatrix}-1&-1\\1&0\end{pmatrix},\quad
 E=SR,\quad T=RS.
                                                        \tag{C.40.7}
\]
Here $T^2=1$, $TST=E$, and $TET=S$. Thus the finite
maps belonging to a sibling pair have the same kernel and
are related by $\operatorname{Int}(T)$. The three distinct
finite maps give the image $Q_0\subset S_3^3$ generated by
\[
             X=(S,R^{-1},E),\qquad Y=(E,S,R^{-1}).
                                                        \tag{C.40.8}
\]
This group is
\[
                  Q_0=C_3^3\rtimes C_2^2,
                       \qquad |Q_0|=108.              \tag{C.40.9}
\]
For example, the two sign vectors are $(1,0,1)$ and
$(1,1,0)$. Their commutator has a nonzero entry in each
of the three $C_3$ coordinates. The three distinct
nontrivial characters of $C_2^2$ on these coordinates
separate them over $\mathbb F_3$, so conjugates of that
commutator span $C_3^3$. This proves (C.40.9).

Let $K_0=\ker(F\to Q_0)$. The six kernels form the
orbit under conjugation by $\Gamma$, hence $K_0$ is
normal in $\Gamma$. Consequently
\[
       \widetilde Q_0=\Gamma/K_0,
                 \qquad |\widetilde Q_0|=648.          \tag{C.40.10}
\]
Only this finite extension of $Q_0$ by the old $S_3$
is asserted. No splitting inside the original modular
group is assumed. For an exact finite verification,
\texttt{pbc\_joint\_local\_descent.py} enumerates the
$108$ triples, verifies every generator edge for the
actions $X\leftrightarrow Y$ and
$(X,Y)\mapsto((XY)^{-1},X)$, and obtains a group of
$648$ actual permutations. Its inner subgroup has order
$108$ and $Q_0$ has trivial centre. Thus this also gives
a faithful explicit model of the finite extension.

\paragraph{The joint local quotient has seven independent free factors.}
Fix $p\geq5$ and put
\[
 E_p=\Gamma'/\ker(K'_2\to K_2'{}^{(p)}),\qquad
 P_p=K_2'{}^{(p)}.
                                                        \tag{C.40.11}
\]
The group $P_p$ is free pro-$p$ of rank two and is the
normal kernel of $E_p\to S_3$. It has a complement
because six is prime to $p$.

For one sibling pair, the two original filling kernels
together impose $x^2=y^2=z^3=1$. This is the finite
triangle group of type $(2,2,3)$, namely $S_3$.
The relative pro-$p$ kernels contain these original
kernels and are contained in the common finite
$S_3$ kernel. Their product is therefore precisely
that finite kernel. Goursat's lemma gives the whole
fibre product $E_p\times_{S_3}E_p$, with the finite
identification specified by $T$ in (C.40.7).
The image of $K_0$ in its pro-$p$ kernel $P_p^2$ is
also onto: the quotient by that image is a quotient
of $F/K_0$, so it is both pro-$p$ and of order prime
to $p$, and hence is trivial.

Let $J_p$ be the image of all six fillings in $E_p^6$.
Its finite image is the actual $Q_0$ just specified.
The image of $K_0$ is a closed normal pro-$p$ subgroup
of $P_p^6$. On the Frattini quotient of this product,
the three sibling blocks are three pairwise
nonisomorphic absolutely irreducible standard
$\mathbb F_p[Q_0]$ modules, each with multiplicity two.
They are distinct because the three $C_3$ coordinates
act on different blocks. Maschke's theorem and the
surjectivity to each whole sibling pair show that the
Frattini image is the entire direct sum. The Frattini
criterion for a finitely generated pro-$p$ group now
gives
\[
                1\longrightarrow P_p^6
                  \longrightarrow J_p
                  \longrightarrow Q_0\longrightarrow1.
                                                        \tag{C.40.12}
\]
Each standard module has zero $Q_0$ coinvariants.
The prime-to-$p$ quotient therefore gives
$H_1(J_p,\mathbb F_p)=0$. In particular $J_p$ has no
nontrivial pro-$p$ quotient: any such quotient would
have a nontrivial finite $p$-group quotient and then
a quotient of order $p$.

The joint image of $F$ in $F^{(p)}\times J_p$ is thus
the full product, again by Goursat. The kernels of
the six fillings form a deck orbit, and the ordinary
pro-$p$ kernel is characteristic. Their intersection
is normal in $\Gamma$. The resulting actual quotient
$Q_p$ of $\Gamma$ has the exact sequence
\[
 1\longrightarrow\Pi_p\longrightarrow Q_p
        \longrightarrow\widetilde Q_0\longrightarrow1,
 \qquad
 \Pi_p=P_{0,p}\times
       \prod_{k=0}^2(P_{k,p}\times P'_{k,p}),            \tag{C.40.13}
\]
where all seven factors are free pro-$p$ groups of
rank two. Here $P_{0,p}$ is the ordinary quotient of
the old $F$. The restriction to $F$ retains its map
to that factor and all six original fillings.

For finitely many distinct primes, the joint image
of $K_0$ in these kernels is the product of the
$\Pi_p$: finite quotients of different prime orders
have no common quotient, and compactness passes from
finite quotients to the pro-prime groups. We therefore
obtain one actual quotient $Q$ with finite base
$\widetilde Q_0$ and normal kernel
\[
                         \Pi=\prod_{p\in\mathcal S}\Pi_p.
                                                        \tag{C.40.14}
\]

\paragraph{The central homology vanishes at the retained primes.}
We compute the continuous homology with compact
$\mathbb Z_p$ coefficients. Each free pro-$p$ factor
has homological dimension one. The completed Kunneth
formula gives
\[
 H_2(\Pi_p,\mathbb Z_p)=
   \bigoplus_{i<j}
       H_1(P_{i,p},\mathbb Z_p)\otimes_{\mathbb Z_p}
       H_1(P_{j,p},\mathbb Z_p).                        \tag{C.40.15}
\]
The tensor factors are finite free lattices; their
ordinary and completed tensor products agree. There
are no single-factor $H_2$ terms and no Tor terms.
Finite-group averaging applies because $648$ is a
unit in $\mathbb Z_p$.

The old $Q_0$ acts trivially on $H_1(P_{0,p})$.
Its three other standard isotypic types are distinct.
Consequently taking $Q_0$ invariants in (C.40.15)
leaves exactly the three sibling invariant lines.
Here is the integral calculation of their actual
deck action. In the bases (C.40.7), put
\[
 B=\begin{pmatrix}2&-1\\-1&2\end{pmatrix},\qquad
 C=BT^{\mathsf t}=
               \begin{pmatrix}-1&-1\\-1&2\end{pmatrix}.
                                                        \tag{C.40.16}
\]
For $A=S,R$ one has $ABA^{\mathsf t}=B$.
The two actions in the sibling pair are $A$ and
$TAT^{-1}$, so its invariant tensor is $C$. It is
primitive and symmetric. These matrices describe
the actual standard representation: in the original
column basis, $\rho$ acts by
$\left(\begin{smallmatrix}-1&1\\-1&0\end{smallmatrix}\right)$;
the basis matrix
$\left(\begin{smallmatrix}1&2\\2&1\end{smallmatrix}\right)$,
of determinant $-3$, changes it to $R$ and leaves $S$
unchanged. It is invertible over every $\mathbb Z_p$
under consideration.

The actual involution $t_k$ in (C.40.6) interchanges
the two sibling coordinates without an additional
linear map. Interchanging two degree-one homology
factors acts on their cross term as minus transpose.
It therefore acts by $-1$ on that invariant line.
Since this pair is fixed as a set, no other line is
sent into it; an invariant vector must have zero
coefficient on that line. Applying this for each
$k$ kills all three lines. Thus the full
$\widetilde Q_0$ invariants, and equivalently its
coinvariants, vanish. Hochschild--Serre gives
\[
                H_2(Q_p,\mathbb Z_p)=0,
                    \qquad H_1(Q_p,\mathbb Z_p)=0.      \tag{C.40.17}
\]
The latter assertion also follows directly from the
standard representations, including the ordinary
factor on which the old $S_3$ acts in its standard
representation.

At a coefficient prime $q\ne p$, a pro-$p$ group
has zero positive compact $\mathbb Z_q$ homology:
the averaging maps on its finite quotients use only
orders invertible in $\mathbb Z_q$, and are compatible
with passage to the limit. For (C.40.14) the same
argument removes every prime except the coefficient
prime. Therefore, for every prime $q$ and $i=1,2$,
\[
 H_i(Q,\mathbb Z_q)\xrightarrow{\sim}
                 H_i(\widetilde Q_0,\mathbb Z_q).       \tag{C.40.18}
\]
At a retained prime both sides are zero; at every
other prime the normal kernel is acyclic.

\paragraph{The entire source-centralization is a pullback.}
Let $N=\ker(\Gamma\to Q)$ and form the two actual
centralizations
\[
 H=\Gamma/\overline{[\Gamma,N]},\quad K=N/\overline{[\Gamma,N]},
 \qquad
 H^0=\Gamma/\overline{[\Gamma,K_0]},\quad
 K^0=K_0/\overline{[\Gamma,K_0]}.
                                                        \tag{C.40.19}
\]
The compact abelian kernel has its usual prime
components $K_q$. As in B.27, the finite cyclic
periodic resolution and the free-product homology
sequence give $H_2(\Gamma,\mathbb Z_q)=0$ for every
$q$. The actual five-term sequence is therefore
\[
 0\longrightarrow H_2(Q,\mathbb Z_q)
  \longrightarrow K_q\longrightarrow H_1(\Gamma,\mathbb Z_q)
  \longrightarrow H_1(Q,\mathbb Z_q)\longrightarrow0.
                                                        \tag{C.40.20}
\]
Use the same sequence for $\widetilde Q_0$ and $K^0$.
Naturality and (C.40.18) show that $K_q\to K_q^0$ is
an isomorphism for every $q$. This is deliberately
the full five-term argument: $N\subset\Gamma'$ is
not assumed, and no identification with a universal
central extension of an arbitrary quotient is used.
Comparing kernels and quotients gives the actual
Cartesian square
\[
                    H\simeq Q\mathbin\times_{\widetilde Q_0}H^0.
                                                        \tag{C.40.21}
\]
In fact $K^0$ is finite and has only two- and
three-primary parts, by (C.40.20) for the finite
$\widetilde Q_0$ and $\Gamma^{\rm ab}=C_6$.

It follows that every automorphism $a$ of $Q$ that
is identity on $\widetilde Q_0$ has a lift
$(a,\operatorname{id})$ to $H$. If its three positive
peripheral conjugators lie in $\Pi$, then their pairs
$(b_s,1),(b_r,1),(b_\tau,1)$ are actual elements of
this fibre product and are peripheral conjugators
for that lifted automorphism. This verifies the
central refinement, including the peripheral data.

We recall precisely why this yields a full lift.
For any finite source quotient $\Gamma\twoheadrightarrow E$,
let $A$ be the image of $N$ and $D=E/A$. In the
finite scalar-cover construction of P.3, every
auxiliary source map to a central extension of $D$
sends $N$ centrally. It thus factors through $H$.
Precomposing by the actual peripheral automorphism
of $H$ supplies the new epimorphism and the three
specified conjugators in every auxiliary cover.
The trace count and subgroup-by-subgroup Mobius
inversion in P.3 apply because they require only
normality of $A$. They give an epimorphism to $E$
with the specified retained values and frames.
Adjoin any finite retained quotient of $Q$ to $E$
when necessary. Compactness of the three conjugator
variables, imposing generation in every finite
quotient, then gives a surjective continuous
endomorphism of the actual modular source. Its
finite generation makes it an automorphism. This
is the modular version of B.11 with arbitrary
specified $N$, and proves a full lift retaining
the action on $Q$ and the three frame values.

\paragraph{Constructing arbitrary prescribed two local actions.}
Now prescribe $\alpha_p,\beta_p\in G_{0,p}$ of
multiplier one for every $p\in\mathcal S$.
An actual local gentle class extends to the relative
modular group $E_p$ with identity quotient action on
$S_3$, by the two cycle formulas of P.4, or by the
normalizer proof of C.1. It preserves its three
positive orbifold peripheral classes. After an
inner correction in $P_p$, it fixes a chosen finite
$S_3$ complement pointwise: two such complements
are conjugate by the normal pro-$p$ group. Its
restriction to $P_p$ therefore commutes with that
complement, and represents the same local gentle
outer class.

Choose a complement $\widetilde Q_0$ in $Q$ to
$\Pi$. Existence and conjugacy follow from the
coprime complement theorem, first at finite
quotients and then by compactness; all primes of
the kernel are at least five. On $P_{0,p}$ its
action factors through the old $S_3$: the kernel
$Q_0$ would otherwise act through a finite
prime-to-$p$ subgroup of the inner pro-$p$
automorphism group, which must be trivial.
Choose the representative of $\alpha_p$ commuting
with this complement action.

The six other factors form a transitive orbit of
$\widetilde Q_0$, with stabilizer $Q_0$. On a
chosen factor its stabilizer action factors through
its actual finite target $S_3$. Indeed the image
of this finite stabilizer in its relative target
$E_p$ is a finite complement; its kernel in $P_p$
is trivial. Choose a representative of $\beta_p$
commuting with that action, and transport it to
the other five factors using the actual complement
action. Stabilizer commutation makes the transport
independent of the chosen transporting element.
All transported maps preserve the corresponding
positive cusp classes; the transport comes from
the six actual marked fillings, up to their inner
target changes.

The product of these seven automorphisms, for all
$p\in\mathcal S$, commutes with $\widetilde Q_0$.
Extending it by identity on this complement gives
an actual automorphism $a$ of $Q$. Its ordinary
local outer coordinate is $\alpha_p$, and its
selected filled coordinate is $\beta_p$.

We check all three required modular conjugators
in the kernel $\Pi$. The elements $a(s)$ and $s$
have order two above the same finite base element,
and $a(r)$ and $r$ have order three. Coprime
complement conjugacy supplies $\Pi$-conjugators.
For the third element put $\tau=sr$, so
$\tau^2=y$. The ordinary coordinate of $a(y)$ is
conjugate to $y$ by a pro-$p$ element because
$\alpha_p$ is peripheral. In each filled coordinate,
the image of $y$ is one of the three orbifold
peripheral elements, and the corresponding extended
$\beta_p$ gives a pro-$p$ conjugator. These independent
coordinate conjugators assemble to an element of
$\Pi$. Thus $a(\tau^2)$ is $\Pi$-conjugate to
$\tau^2$.

After this inner adjustment, suppose
$a(\tau^2)=\tau^2$. Write $a(\tau)=v\tau$.
Then $v\in C_\Pi(\tau^2)$, and
\[
                  v\operatorname{Int}(\tau)(v)=1.
                                                        \tag{C.40.22}
\]
The automorphism $\operatorname{Int}(\tau)$ has
order dividing two on this centralizer, since
its square is conjugation by $\tau^2$.
Coprime complement conjugacy gives
$H^1(C_2,C_\Pi(\tau^2))=1$; it can equivalently
be applied in each pro-$p$ factor and assembled.
Equation (C.40.22) is therefore a coboundary,
and an additional $\Pi$-conjugator carries
$a(\tau)$ to $\tau$. This proves positive peripheral
conjugacy of the third element with a pure retained
frame. The same argument, with one target factor,
justifies the pro-$p$ frame for its infinite
orbifold cusp used in the preceding paragraph.

By (C.40.21), these three frames lift as
$(b_i,1)$ to the entire source-centralization.
The finite refinement and compactness argument
therefore gives an actual positive peripheral
automorphism $\Phi$ of the full $\Gamma$, retaining
$a$ and those frames. Its quotient action on the
old $S_3$ is identity. Restricting to $F$ gives
a pure peripheral outer class centralizing the
geometric $S_3$, hence a full gentle element of
multiplier one. Equivalently, fixing the modular
cusp by an inner correction gives the exact full
two- and three-cycle equations as in P.4; standard
normalization gives the conjugator in $F'$.
These are full equations, not only equations in
the retained local quotients.

The ordinary retained quotient is $F^{(p)}$, so
the first local outer class is exactly $\alpha_p$.
The selected filling is the original $\kappa$;
its relative target action restricts to
$\beta_p$ on its characteristic pure subgroup.
This is precisely the operation $D$ of C.1.
Inner normalization of the full representative
does not change either outer coordinate. We have
therefore realized every prescribed finite tuple
$(\alpha_p,\beta_p)_{p\in\mathcal S}$.

\paragraph{All primes and all common multipliers.}
The full gentle group is compact, and the local
coordinate maps and $D$ are continuous. Surjectivity
for every finite set of primes therefore gives
surjectivity onto the whole multiplier-one product.
For a tuple of arbitrary common multipliers
$(\lambda_p)_{p\geq5}$, choose a global unit with
those components and an actual Galois element
having that cyclotomic character. The cyclotomic
character is onto. Its two coordinates agree
because $D$ fixes the specified Galois image.
Divide each desired pair by these common Galois
coordinates, apply the multiplier-one result,
then multiply by the same Galois element.
This proves (C.40.2). The reverse inclusion is
the multiplier preservation already proved in C.1.
\end{proof}
